\subsection[The propagated pair (2,3)]{The propagated pair $(2,3)$}\label{loc:S23}
In this proof $\tau=d-\delta_{2}$.
Here $P=h_2f_3$. We prove $\E[p(T^LP)^4]\le1$ for every $L\ge1$,
with the same actual directions at every length and parameter.
Put $m=1-d/2$, $k=1/m$ and $s=K^2g$. Since
$h_2+f_3=2-d+s$, positivity and AM--GM give
\[
 P\le F_0=(m+s/2)^2,\qquad Y_L=T^LF_0\ge T^LP\ge0.
\]
We shall compare all moments with
\begin{equation}\label{eq:compact-23-taylor}
 m^{-8}\ge1+4d+9d^2,\qquad dm^8<1/9.
\end{equation}
The first is the positive Taylor remainder of $(1-d/2)^{-8}$.
The second follows by differentiation: its maximum is
$(2/9)(8/9)^8<1/9$ at $d=2/9$.

\proofpart{Finite spectra and every $L\ge4$}
Let $Q_n(z)=G_n(z)-kz^{n+2}$. For every $n\ge3$ its spectral norm
is at most one. To see this write
\[
 Q_n=\frac{1-z^n(1+kz^2-kz^3)}{1-z},\quad
 A_x=1-kx\ge41/80,\quad C_x=A_x^2+k^2x^3\le1.
\]
For $n\equiv0,3\pmod4$ the numerator of the squared modulus is at most
one. For residues $1,2$ the losses after division by $x$ are at least
$1-(12/5)x^3-x^4\ge51/80$ and
$1-2x^2-x^5\ge15/32$, respectively. This includes every finite length.
For $V_L=S'_L-kK'^{L+3}\one$, $L\ge3$, finite alternating means give
$\alpha=1-\E V_L,\eta=\E[g'V_L]\in[\tau,d]$.
Even added powers are $\beta_{2j-2}+k(-1)^j\delta_{j}$; pair consecutive terms,
using $k\theta^2<1$, with the first case $d-k\delta_{3}\ge\tau$.
Hence $\|V_L-1\|^2\le d$ and, for $L\ge4$,
$|V_L-q|\le\sqrt d$.

Duhamel and $T(s^2)\le2\delta_{3}$ give
\[
 Y_L\le m^2V_L-m\delta_{2}\bar f_{L-2}+(\delta_{3}/2)\bar f_{L-1}
                 -\sum_{j=3}^L\omega_j\bar f_{L-j},
\]
where $\omega_j=m^2\beta_j$ for odd $j$, and
$\omega_j=m^2\beta_j-m(-1)^{j/2}\delta_{j/2+1}$ for even $j$.
All are at least $m^2\tau$: for $j\equiv0\pmod4$ pair the final
positive beta term with $\delta_{j/2+1}$, using $m>\theta^2$;
the boundary $j=4$ uses $m>\theta$.
For $L\ge4$ this implies $Y_L\le m^2Z$ with
\[
 Z=V_L-a-\tau q,\quad a=\tau-\delta_{3}/m^2\ge2279d/6400>d/3.
\]
Then $\mu=\E Z\le1-d$, covariance $\E[g'(Z-\mu)]\ge0$,
$\operatorname{Var}Z\le d$, and
$Z\le(1-d/3)q-d/3+\sqrt d$.
Profile A of Section~\ref{calc:coefficientprofiles} bounds the fourth
coefficient by seven for exactly this mean and row envelope.
Thus $\E[pY_L^4]\le m^8[(1-d)^4+7d]\le1$;
the Taylor difference is $d+3d^2+4d^3-d^4\ge0$.

\proofpart{Three and two steps}
For each short case below define $H=p(B^2+2uB+3u^2)$ using that case's stated mean upper bound $u$ and pointwise upper bound $B$.
For $L=3$, retain the actual term $-m\delta_{2}q$:
\[
 Y_3\le m^2Z,\quad Z=V_3-a-bq,\quad
 a=d-\delta_{3}/m^2\ge4879d/6400,\quad b=k\delta_{2}.
\]
The mean is at most $1-d$, covariance at least $3d/10$, and variance
at most $d$. The tail $Q_2$ has squared modulus
$(1-kx^2)^2+x\le233/200<(13/12)^2$.
Thus $Z\le q-3d/4+13\sqrt d/12$. The coefficient certificate is
\begin{align*}
 D&=-1/4-407x^2/144-91x^3/24+81x^4/16,\\
 Q&=-5/8+39x/8-1163x^2/96-91x^3/16+243x^4/32.
\end{align*}
Here $D=H'(3/2)$, $Q=H(3/2)-7$ and $H''\le-3$ on the maximizing
interval. For $x\le1/3$, $-D<2/3$, $Q<-1/8$, giving margin $-11/216$.
For $x\ge1/3$, convexity and negative endpoints of $Q'$ give
$Q\le-25/54$, $-D<7/5$, giving margin $-92/675$.
The same Taylor comparison proves the three-step case.

For $L=2$, set $\varepsilon=\delta_{3}/(2m^2)$ and
$Z=V_2-k\delta_{2}+\varepsilon q\ge Y_2/m^2$.
Then
\[
 \mu\le1-d,\qquad h\ge4d/5,\qquad
 \operatorname{Var}Z\le25d/22,\qquad Z\le q+21\sqrt d/20.
\]
For the variance, use $v=\delta_{2}/d$ and
\[
 d+\delta_{2}-2k\delta_{3}+k^2\delta_{5}\le d[1+v-(2-\theta^2)v^2]<25d/22;
\]
the remaining $-d^2+2\varepsilon(d-k\delta_{3})+\varepsilon^2d$ is
nonpositive since $\varepsilon<3d/25$.
The row bound is $\sqrt{d+k^2\delta_{4}}\le21\sqrt d/20$.
For $\lambda=21/20$ the corresponding coefficient function is concave,
has maximum below $3/2$ and curvature at most $-3$; its derivative there
has absolute value below $4/3$, while its value is at most $57/8$.
The tangent bound $57/8+8/27<15/2$ proves $M\le15/2$.
The endpoint comparison is therefore
\[
 \E[pY_2^4]\le m^8\max\{Cd,(1-d)^3(1-21d/5)+Cd\},\quad C=375/44.
\]
The first branch is below $125/132$ by \eqref{eq:compact-23-taylor}.
For the second, its difference below $1+4d+9d^2$ is at least
$21d/44+61d^3/5\ge0$.

\proofpart{One step}
Exactly $Y_1=m^2Z$, where
\[
 Z=q-kK^3g+e,\quad e=(\delta_{3}-K(s^2))/(4m^2),\quad
 0\le e\le \delta_{3}/(2m^2).
\]
The mean, covariance and centered variance obey
\[
 \mu\le1-3\delta_{2}/4,\quad h\ge d-\delta_{2}/10,\quad
 \operatorname{Var}Z\le d+7\delta_{2}/16.
\]
Indeed the mean loss and covariance error are bounded by
\[
 (1-\theta/2)\delta_{2}=51\delta_{2}/64,\qquad
 (117/800)\sqrt d\,\delta_{2}<\delta_{2}/10,
\]
respectively.
For the variance, use the square of
\[
 \sqrt{d+k^2\delta_{4}}+\delta_{3}/(4m^2).
\]
The three added coefficients are
$1521/6400$, $2457\sqrt d/8000$, $13689\delta_{2}/640000$,
whose sum is below $751/1750<7/16$.
Put $u=\sqrt{\delta_{2}}<3/8$. The row envelope is
$Z\le q+4u/5+3u^2/10$.
For this envelope and mean $1-3u^2/4$, the concave fourth coefficient
has curvature at most $-3$ below $3/2$ and
\begin{align*}
 D&=-1/4-59u^2/25-18u^3/25+531u^4/400,\\
 H(3/2)&=51/8+18u/5-1113u^2/200-27u^3/25+1593u^4/800.
\end{align*}
$-D<5/8$; the cubic and quartic sum is nonpositive, and the quadratic
gain is below $3/5$. Hence $M<51/8+3/5+25/384<57/8$.
The zero mean endpoint is below $247/256$ by $\delta_{2}\le d/2$ and
\eqref{eq:compact-23-taylor}.
For the other endpoint, expand with $\delta=3\delta_{2}/4$:
\[
 (1-\delta)^3(1-\delta-4h)
 \le1-4d-13\delta_{2}/5+9d\delta_{2}+99\delta_{2}^2/40.
\]
Adding $M(d+7\delta_{2}/16)$ and subtracting from the Taylor lower bound
leaves at least $789d/1280+621d^2/160\ge0$.
This completes every positive length. Equal pairs are restored by
positivity, full reversal exchanges $p,T$ with $q,A$, and $d=0$ gives equality.
